\begin{abstract}
Model merging has emerged as a powerful paradigm for combining the specialized capabilities of pre-trained models. However, standard dynamic ensembling methods have become increasingly convoluted, relying on over-parameterized routing networks that require expensive calibration datasets, backpropagation optimization, and complex training heuristics. In this work, we apply a relentless application of Occam's razor to dynamic model ensembling: \textbf{Parameter-Free Task-Space Projection (PFSR)}. PFSR is a completely training-free, data-free, closed-form linear projection that extracts task-representative centroids from frozen experts and projects online feature representations to derive ensembling coordinates in a single forward pass. To investigate whether orthogonalizing these task coordinates can eliminate cross-talk under task overlap, we analyze an advanced extension using \textbf{L{\"o}wdin Symmetric Orthogonalization (OTSP)}. We prove mathematically and empirically that under symmetric task correlations, OTSP and PFSR have identical Signal-to-Noise Ratios (SNR) and routing decisions, making orthogonalization redundant. In asymmetric task environments, we show that OTSP systematically underperforms PFSR by 0.2\% to 1.6\% due to the \textit{Noise Amplification Penalty} and \textit{Noise Spillover Penalty} of the orthogonalized basis under active representation noise. Crucially, we transparently deconstruct the limitations of disjoint orthogonal sandboxes, demonstrating how the \textit{Orthogonal Masking Effect} flattens joint classification performance across all normalized routers (matching the expert ceiling of 74.46\%), and we reframe the unconstrained linear router's instability as a pedagogical demonstration of the necessity of the probability-simplex constraint for vectorized online deployment ($B=1$). In asymmetric environments, we show that while static Uniform Merging slightly dominates average classification accuracy under high overlap due to ensembling redundancy, PFSR and OTSP act as robust, zero-parameter mechanisms for routing specificity and noise isolation. Our results demonstrate that the simpler, raw unorthogonalized projection is not only computationally cheaper, but systematically more robust than more complex orthogonalized variants, prompting a critical re-evaluation of complexity in model-merging architectures.
\end{abstract}
